% Generated draft: requires mathematical and source review.
\documentclass[11pt]{article}
\usepackage[margin=25mm]{geometry}
\usepackage{fontspec}
\setmainfont{Latin Modern Roman}
\usepackage{amsmath,amssymb,mathtools,unicode-math}
\setmathfont{Latin Modern Math}
\usepackage{enumitem,needspace,xurl,hyperref}
\setlength{\parindent}{0pt}
\setlength{\parskip}{5pt}
\setlength{\emergencystretch}{3em}
\newcommand{\N}{\mathbb{N}}
\newcommand{\Z}{\mathbb{Z}}
\newcommand{\Q}{\mathbb{Q}}
\newcommand{\R}{\mathbb{R}}
\newcommand{\C}{\mathbb{C}}
\newcommand{\F}{\mathbb{F}}
\newcommand{\A}{\mathbb{A}}
\newcommand{\PP}{\mathbb{P}}
\newcommand{\E}{\mathbb{E}}
\newcommand{\eps}{\varepsilon}
\newcommand{\dd}{\,\mathrm{d}}
\begin{document}
\section*{UnsolvedMath ID: 800037}
\textbf{Source problem code:} catalogue entry
\subsection{Definitions and mathematical statement}
A semigroup is a nonempty set equipped with an associative binary operation. For a class \(\mathcal K\) of semigroups, let \(\mathsf S(\mathcal K)\) be the class of semigroups isomorphic to nonempty subsemigroups of members of \(\mathcal K\), and let \(\mathsf H(\mathcal K)\) be the class of semigroups that are surjective homomorphic images of members of \(\mathcal K\). For a class operator \(\mathsf O\), closure under \(\mathsf O\) means \(\mathsf O(\mathcal K)\subseteq\mathcal K\).

\textbf{Problem (G. Mashevitzky).} Find a characterization of inclusive varieties of semigroups analogous to Birkhoff's characterization of varieties by closure operators. Specifically, determine a set \(\mathscr O\) of Birkhoff type operators such that, for every class \(\mathcal K\) of semigroups,
\[
\mathcal K\text{ is an inclusive variety}
\quad\Longleftrightarrow\quad
\mathsf O(\mathcal K)\subseteq\mathcal K
\text{ for every }\mathsf O\in\mathscr O.
\]
The source additionally asks, parenthetically, for a minimal such set of operators. It reports that inclusive varieties are closed under subsemigroups and homomorphic images; thus it supplies the necessary conditions
\[
\mathsf S(\mathcal K)\subseteq\mathcal K,
\qquad
\mathsf H(\mathcal K)\subseteq\mathcal K
\]
when \(\mathcal K\) is an inclusive variety.

\textbf{Unresolved qualifications.} The supplied record does not define an inclusive variety. Consequently, the displayed equivalence is a schematic statement of the requested characterization, not a complete intrinsic formulation of its left-hand side. The record also does not specify the permitted Birkhoff type operators or the meaning of minimality. These specifications must be recovered from an authoritative definition or fuller source context before a precise closure characterization can be formulated.
\subsection{Short English statement}
Find closure operators that characterize inclusive varieties of semigroups, in analogy with Birkhoff's theorem, and seek a minimal collection of such operators. The source reports closure under subsemigroups and homomorphic images. A complete formulation remains blocked by the missing definition of inclusive variety and the unspecified meanings of Birkhoff type operators and minimality.
\subsection{Sources}
\begin{itemize}
\item Local UnsolvedMath registry, permanent ID \(800037\), exact source problem code \texttt{catalogue entry}, title \emph{The Lyapin's Notebook --- Problem 39}; collection: AMR Open Problem Lists. The supplied registry labels the problem open; that label is supplied curation, not an independently verified mathematical or literature assessment.
\item Bershadsky, Kublanovsky, Mashevitzky, \emph{The Lyapin's Notebook: Unsolved Problems in Semigroup Theory} (2026), Problem 39. Supplied source URL: \url{https://arxiv.org/abs/2604.04763}. The problem statement attributes the problem specifically to G. Mashevitzky; the registry's proposed-by field lists Bershadsky, Kublanovsky, Mashevitzky and gives proposed year 2026.
\item Supplied provenance records access on 2026-10-16 and extraction by \texttt{arxiv\_latex}. This access date and provenance are supplied source metadata; no independent source visit or verification is claimed.
\end{itemize}
\end{document}
